\documentclass[a4paper,10pt]{article}

\usepackage[utf8]{inputenc}
\usepackage[T1]{fontenc}
\usepackage[english]{babel}
\usepackage[margin=0.85in]{geometry}
\usepackage{titlesec}
\usepackage{enumitem}
\usepackage{hyperref}
\usepackage{xcolor}
\usepackage{indentfirst}

\titleformat{\section}{\color{blue}\large\bfseries}{}{0em}{}[\titlerule]
\pagenumbering{gobble}
\setlist[itemize]{leftmargin=1.5em, itemsep=0.15em, topsep=0.2em}

\begin{document}

\begin{center}
    {\LARGE \textbf{Pedro Henrique Klein}}\\
    \vspace{0.2cm}
    \href{mailto:pedrohenriquekleinphg@gmail.com}{pedrohenriquekleinphg@gmail.com} \textbar\
    (54) 99704-9008 \textbar\
    \href{https://www.linkedin.com/in/pedro-henrique-klein-a41122221/}{LinkedIn} \textbar\
    \href{https://github.com/Pedrinhonitz}{GitHub} \textbar\
    Chapecó, SC, Brazil
\end{center}

\vspace{0.3cm}

\section*{About}

I started programming at 15, first with Arduino and then with Python on my own. At 16 I joined the market as a developer and, until 18, I worked on web backends with Django, Node.js, PHP, JavaScript, and TypeScript.

In 2021 I started Computer Science at Universidade Federal da Fronteira Sul (UFFS), in Chapecó, where I still study. In early 2022 I left backend work for data engineering. Since then I have been building pipelines, modeling data, and orchestrating processes with Python, SQL, Airflow, GCP, and AWS.

That mix of software and data lets me move across technical areas and work well with different teams.

\section*{Education}

\textbf{Universidade Federal da Fronteira Sul (UFFS)} \hfill 2021 to 2027 \\
Bachelor's degree in Computer Science, Chapecó campus

\section*{Professional experience}

\textbf{BIP Brasil, Pro Carbono project (Bayer)} \hfill Sep 2024 to Aug 2026

\textit{Senior Data Engineer} \hfill Aug 2025 to Aug 2026
\begin{itemize}
    \item I build and maintain pipelines with Apache Airflow, PySpark, AWS Glue, and Astronomer Cosmos.
    \item I create internal applications and APIs with Python, FastAPI, and TypeScript to automate processes and integrate systems.
    \item I model data with DBT, with quality, structure, and lineage of transformations.
    \item I administer and tune Linux, Docker, and Kubernetes environments for data engineering workloads.
    \item I write and tune queries in PostgreSQL, Amazon Athena, and BigQuery for large datasets.
    \item I set up CI/CD pipelines, review code, and define engineering practices.
    \item I contribute to Apache Airflow with fixes and improvements for the community.
\end{itemize}

\textit{Pleno Data Engineer} \hfill Sep 2024 to Aug 2025
\begin{itemize}
    \item I implemented and maintained pipelines with Apache Airflow and PySpark.
    \item I built processing and storage solutions on AWS and Google Cloud Platform.
    \item I modeled data with DBT and Elementary.
    \item I shipped CI/CD with GitHub Actions, Jenkins, and Terraform, plus coding standards and pull request reviews.
\end{itemize}

\textbf{Descomplica} \hfill Jan 2022 to Sep 2024 \\
\textit{Pleno Data Engineer}
\begin{itemize}
    \item I orchestrate Apache Airflow workflows to run and monitor data pipelines.
    \item I package workloads with Docker and orchestrate them on Kubernetes.
    \item I administer and tune Linux environments.
    \item I build data warehouse solutions with BigQuery, Amazon Athena, and Snowflake.
    \item I create dashboards and reports in Metabase.
    \item I automate ETL with Python and Shell.
    \item I write and tune SQL for large datasets.
\end{itemize}

\textbf{IXC-Soft} \hfill Nov 2021 to Jan 2022 \\
\textit{Software Engineering Intern}
\begin{itemize}
    \item I built backend features in PHP.
    \item I built interfaces with Twig.
    \item I packaged development environments with Docker.
    \item I administered and tuned MariaDB databases.
\end{itemize}

\section*{Open Source}

\textbf{\href{https://github.com/apache/airflow/pull/66431}{Apache Airflow (Apache Software Foundation)}} \\
Pull request merged into the official repository.
\begin{itemize}
    \item I fixed the scheduler handling of \textit{deferred tasks}, removing a race condition with stale events from the executor.
    \item I worked with maintainers through technical review until the change was approved and merged into the main codebase.
\end{itemize}

\section*{Technical skills}

\begin{itemize}
    \item \textbf{Languages:} Python, SQL, TypeScript, JavaScript, Bash, Go, Java, C, C\texttt{++}, Haskell, and Ruby.
    \item \textbf{Data engineering:} Apache Airflow, Astronomer Cosmos, DBT, Elementary, PySpark, AWS Glue, BigQuery, Snowflake, Amazon Athena, Luigi, and Dagster.
    \item \textbf{DevOps and infrastructure:} Docker, Kubernetes, Terraform, GitHub Actions, Jenkins, and Linux.
    \item \textbf{Databases:} PostgreSQL, SQL Server, Oracle, and Neo4j.
    \item \textbf{Cloud:} Amazon Web Services (AWS) and Google Cloud Platform (GCP).
\end{itemize}

\section*{Languages}

\begin{itemize}
    \item Portuguese: native
    \item English: A2
    \item Spanish: B1
\end{itemize}

\end{document}
